%% CPP 2027 Submission — 0pat Exercise XXX: Transposition and the Origin in the Orthogonal Euler Circle Frame
\documentclass[sigplan,10pt,anonymous,review]{acmart}
\settopmatter{printfolios=true,printacmref=false}

\AtBeginDocument{%
  \providecommand\BibTeX{{\normalfont B\scshape ib\TeX}}}

\settopmatter{printacmref=false}
\renewcommand\footnotetextcopyrightpermission[1]{}
\pagestyle{plain}

\usepackage{microtype}
\usepackage{booktabs}
\usepackage{amsmath}
% XeTeX (tectonic): acmart loads newtxmath/libertine which already defines \Bbbk;
% reset it so amssymb's definition is accepted.
\let\Bbbk\relax
\usepackage{amssymb}

\begin{document}

\title{Transposition and the Origin in the Orthogonal Euler Circle\\
Frame (Exercise XXX)}

\author{Anonymous}
\affiliation{%
  \institution{Anonymous}
}
\email{anonymous@example.org}

\begin{abstract}
We formalize the transposition (mirror) structure of the orthogonal
Euler circle frame --- the radius-1 and radius-$i$ circles of the
split-complex numbers (Exercise XXVII) --- and determine the origin of
the frame. Seven theorems, machine-checked with zero \texttt{sorry}:
the transposition $(a,b) \mapsto (b,a)$ is an involution; it maps $1 =
(1,0)$ to $i = (0,1)$; it maps the radius-1 circle $a^2 - b^2 = 1$ to
the radius-$i$ circle $a^2 - b^2 = -1$ in both directions; and it fixes
the origin $(0,0)$. The natural-number lattice of each circle is a
single point: the radius-1 circle contains exactly $(1,0)$ and the
radius-$i$ circle exactly $(0,1)$ --- the two units are transposes of
each other, and the earlier observation that lattice observation
degenerates to $\{1, i\}$ is made precise. The origin of the frame is
$(0,0)$: the fixed point of the transposition and the common center of
the two conjugate hyperbolas. It is neither $1$ nor $i$.
\end{abstract}

\keywords{Lean, formalization, split-complex, transposition, origin, 0pat}

\maketitle

\section{The frame and the question}

The orthogonal Euler circle frame (Exercise XXVII) consists of the
radius-1 circle $a^2 - b^2 = 1$ and the radius-$i$ circle
$a^2 - b^2 = -1$ in the split-complex plane $a + bj$, $j^2 = +1$. The
frame was built to observe the prime axis through two orthogonal
circles with different units ($1$ and $i$). This exercise formalizes the
transposition structure of the frame and answers a positional question:
in this frame, where is the origin?

\section{The theorems}

\paragraph{Transposition (transpose\_involutive, transpose\_maps\_one\_to\_i).}
The map $(a,b) \mapsto (b,a)$ exchanges the coordinates. It is an
involution (applying it twice returns the point), and it maps the unit
of the radius-1 circle to the unit of the radius-$i$ circle:
$1 = (1,0) \mapsto (0,1) = i$.

\paragraph{The circles are transposes (transpose\_maps\_unit\_to\_imaginary,
transpose\_maps\_imaginary\_to\_unit).} A point $(a,b)$ lies on the
radius-1 circle exactly when its transpose $(b,a)$ lies on the
radius-$i$ circle, in both directions:
$a^2 - b^2 = 1 \leftrightarrow b^2 - a^2 = -1$. The two circles are
mirror images under the transposition.

\paragraph{Lattice uniqueness (lattice\_unique\_unit,
lattice\_unique\_imaginary).} Under the constraint that both
coordinates are natural numbers, each circle contains exactly one
point: the radius-1 circle contains only $(1,0) = 1$, and the radius-$i$
circle contains only $(0,1) = i$. The proofs are elementary: for
$(a,b) \in \mathbb{N}^2$, $a^2 - b^2 = 1$ forces $b = 0$ (otherwise the
difference is at least $3$) and then $a = 1$; symmetrically for
$b^2 - a^2 = 1$. This makes precise the earlier observation that
lattice observation of the frame degenerates to the two units
$\{1, i\}$: each circle contributes exactly its own unit, and the two
units are transposes of each other.

\paragraph{The origin (transpose\_origin\_fixed).} The transposition
fixes the origin: $(0,0) \mapsto (0,0)$. The origin is the common
center of the two conjugate hyperbolas $a^2 - b^2 = \pm 1$ --- the
point about which the two circles are symmetric under the transposition
$(a,b) \leftrightarrow (b,a)$ combined with negation. It is neither $1$
nor $i$: those are the units of the two circles, the unique lattice
points, and the transposition pairs them; the origin is their fixed
point.

\section{Honest boundary and position}

The content is elementary algebra (coordinate transposition, lattice
equations in $\mathbb{N}$); the value is positional: the orthogonal
Euler circle frame now has its transposition structure formalized and
its origin determined. The origin of the frame is $(0,0)$ --- the split
zero and the common center of the two circles. The result concerns the
frame, not the Riemann zeta function: no claim is made about zeros, and
the analysis of $\zeta$ remains the open span of the bridge.

\section{Formalization}

The proof is a single Lean file (\texttt{proof.lean}, 185 lines),
importing only \texttt{Mathlib.Data.Real.Basic} and basic tactics. The
split-complex structure (30 lines, as in Exercises XXVII--XXVIII) is
self-contained. Seven theorems, compiled with Lean 4.32.2 / mathlib
v4.32.2, zero errors, zero \texttt{sorry}.

\begin{center}
\begin{tabular}{lll}
\toprule
Theorem & Statement & Proof core \\
\midrule
\texttt{transpose\_involutive} & $T(T z) = z$ & ext; simp \\
\texttt{transpose\_maps\_one\_to\_i} & $T(1) = i$ & ext; simp \\
\texttt{transpose\_origin\_fixed} & $T(0) = 0$ & ext; simp \\
\texttt{transpose\_maps\_unit\_to\_imaginary} & $\|z\|^2 = 1 \Rightarrow \|Tz\|^2 = -1$ & nlinarith \\
\texttt{transpose\_maps\_imaginary\_to\_unit} & $\|z\|^2 = -1 \Rightarrow \|Tz\|^2 = 1$ & nlinarith \\
\texttt{lattice\_unique\_unit} & $a^2 - b^2 = 1 \Rightarrow (a,b) = (1,0)$ & case split; omega \\
\texttt{lattice\_unique\_imaginary} & $b^2 - a^2 = 1 \Rightarrow (a,b) = (0,1)$ & case split; omega \\
\bottomrule
\end{tabular}
\end{center}

\section{Conclusion}

The orthogonal Euler circle frame is now fully determined: its
transposition structure is formalized (involution, $1 \leftrightarrow
i$, circles as mirrors), its lattice observation degenerates to the two
transposed units $\{1, i\}$, and its origin is $(0,0)$ --- the common
center of the two circles, fixed by the transposition, neither $1$ nor
$i$. The frame is complete; the analysis of $\zeta$ remains the open
span of the bridge.

\bibliographystyle{ACM-Reference-Format}
\bibliography{refs}

\end{document}
